\newpage
\chapter{公式資料、規則改正及び公示}

\section{案内を作成するとき}

担当者は、資料に書く制度と運用について、現行の学園憲章、学園規則及び公示済みの案内を確認する。
検討中の人数や次の学期から予定する制度を、既に決定・実施したものとして書かない。
資料には対象、更新日及び参照先を示し、公開前に別の教務主事と条文参照、個人情報の有無及び表現を確認する。
確認は資料の正確さを確かめるためであり、規則にない認可手続を加えるものではない。

\section{規則又はガイドラインを改めるとき}

規則の改正は学園規則第９章に従う。
教務主事の発議、他の教務主事への事前通知による効力発生、学園長による速やかな学生への公示を区別する。
必要に応じて学生の意見を募る。
ガイドラインは第１章第５条に従い、憲章と規則に反せず、規則にない学生の義務を作らない。
学生の活動に関わるガイドラインを制定、改正又は廃止したときは速やかに公示する。
ソースの編集や案の共有を、これらの手続や公示が済んだこととして扱わない。

\section{変更後の確認と記録}

担当者は変更点、根拠、事前通知と公示の時点、改めた関連案内を記録する。
案内を変更するときは、申請フォーム、公式サイト、学内の固定投稿等に古い説明が残っていないか確認し、担当者に修正を引き継ぐ。
変更の理由と影響を受ける学生への連絡を明確にし、未決定の事項は未決定のまま示す。
